% Pattern: kind/untitled. An environment of the plain style with no title and
% a counter of its own is a statement, with a proof or without one. An
% untitled environment that shares the counter of the theorems is a numbered
% paragraph, and so is one whose body is blank. The proof after a statement
% proves it with no flag. Seen in the gold pilot (set X).
% want point-1 statement 1 status=none
% want point-2 statement 2 status=here
% want proof-1 proof - proves=point-2 flags=
% want point-3 unknown 3
% want theorem-1 statement/theorem 1 status=none
% want para-1 unknown 2
\documentclass{article}
\usepackage{amsthm}
\theoremstyle{plain}
\newtheorem{theorem}{Theorem}
\newtheorem{point}{}
\newtheorem{para}[theorem]{}
\begin{document}
\begin{point} A set is finite when it has a bijection to some $n$. \end{point}
\begin{point} Every finite set is countable. \end{point}
\begin{proof} Take the bijection. \end{proof}
\begin{point} \ \end{point}
\begin{theorem} Every countable set is small. \end{theorem}
\begin{para} We now turn to the sets that are not countable. \end{para}
\end{document}
